% SPDX-FileCopyrightText: 2016-2026 Lari Natri <lari.natri@iki.fi>
% SPDX-License-Identifier: GPL-3.0-or-later
%
% Shared language sections, passages and in-song prose environments.
% This file is part of the 'ulsbs' package.
%
% Load: ulsbs.sty, after song helpers and before media in the song-content phase.
% Requires: final typography/style, book policy, Polyglossia, xstring, xparse,
% xcolor, ragged2e, enumitem and verbatim. Layout/verse owners install block
% starts, fill state and verse machinery later, before document-time use.
% State: owns language-section defaults, shared content indent, translation flag
% and prose counters reset by the song lifecycle. Inherits @ as a letter with
% no enclosing group; preserves environment-local groups and raw source skips.
% Passage em lengths are measured now in the post-typography font context;
% change allocated lengths afterward with setlength. Public providecommand
% defaults can be defined before loading or renewed afterward and expand at
% environment entry. Hiding policy and language lists are checked at use time.

\ProvidesFile{ulsbs-prose.code.tex}[ULSBS language sections and song prose]

% === Shared language-section environment ===

% Shared content indentation. Override after loading with
% \setlength{\ulContentIndentWidth}{<dimension>}; used by passages and repeats.
\newlength\ulContentIndentWidth%
\setlength\ulContentIndentWidth{7.33pt}

% Private language-section defaults
% NOTE: These are meant to be locally renewed in environments using the
% ulsbs@langsect environment. Be sure to renew locally only, as other environments
% might depend on these defaults.
\newif\ifulsbs@upcaselangsecttitle\ulsbs@upcaselangsecttitlefalse%
\newcommand{\ulsbs@langsectfont}{}% style for the default text
\newcommand{\ulsbs@langsecttitlefont}{}% style for the title text
% Amount of indentation for default text inside langsect environment
\newlength\ulsbs@langsectleftskipadd\setlength\ulsbs@langsectleftskipadd{0pt}%
\newlength\ulsbs@langsectrightskip\setlength\ulsbs@langsectrightskip{0pt}%
\newcommand{\ulsbs@langsectbaselinemultiplier}{1}%
\newcommand{\ulsbs@langsectbeginextra}{}%
\newcommand{\ulsbs@langsectendextra}{}%
% Omit the displaying of the language codes in the title line for the languages
% listed in this comma separated list. Only the *codes* are omitted, the rest
% of the possible title nor the content is affected.
\newcommand{\ulsbs@langsectomitlangcodes}{}%

% Starts a new subsubsection and prints it's title given as the only parameter.
% Used in the definition of passage environment.
\newcommand\ulsbs@langsectstart{\@startsection{subsubsection}{3}{\ulContentIndentWidth}%\z@
                                 {-3.25ex\@plus -1ex \@minus -.2ex}%
                                 {1.3ex \@plus .2ex}%
                                 {\ulsbs@langsecttitlefont}}



% Returns the formatted version of a language code given as an argument, plus
% a small horizontal space after it. If the code is in the omit list, nothing
% is returned.
\NewDocumentCommand{\ulsbs@langcodeformatted}{m}{%
  \IfValueT{#1}{% must not be empty for string check to work
    \IfSubStr{\ulsbs@langsectomitlangcodes}{#1}{% is in the omit list, return nothing
    }{% current langcode is not in the omit list
      {%
        \colorlet{ulsbs@currcolor}{.}%
        \ulsbs@langcodebox{\color{ulsbs@currcolor}\normalfont\bfseries #1}%
        \hspace{.618034em}%
      }%
    }%
  }%
}



\newcommand{\ulsbs@langcodewithtitle}[2]{%
  \ulsbs@langcodeformatted{#1}#2%
}

% Environment for text that can be put inside a song environment and elsewhere,
% with support for languages. Supports colouring and hyphenation of language
% texts. Intended for use within this package.
% Requires packages: xstring, xparse, xcolor
% Usage: \begin{langsect}*[LANGCODE]{<title>}*
%   - 1st star: language code is displayed (and section started) even if title
%     is empty
%   - If LANGCODE is present, it is displayed before the title in brackets.
%     If it is one of the supported ones: PT, EN, or FI (must be uppercase),
%     the langsect text is coloured and hyphenation for that language is used.
%   - 2nd star: obey lines
%   - If <title> is empty and 1st star is present, no title is displayed
\NewDocumentEnvironment{ulsbs@langsect}{s o m s}{% pre env
  \bgroup%
  \let\ulsbs@oldobeylines\obeylines%
  \let\ulsbs@olddisobeylines\disobeylines%
  \renewcommand{\obeylines}{\ulsbs@oldobeylines\parskip=0pt}%
  \renewcommand{\disobeylines}{%
    \ulsbs@olddisobeylines%
    \justifying% ragged2e package: force justifying
    \parskip=.381966\baselineskip%
  }%
  \newcommand{\polpar}{\par\vspace{.618034\baselineskip}}%
  \renewcommand{\up}[1]{\textsuperscript{\color{ulLanguageSectionMarkColor}##1}}%
  \def\ulsbs@langsecttitle{#3}%
  \IfValueTF{#2}{% optional lang id (must be upcase) signifying translation
    \def\ulsbs@langsectlang{#2}%
    \def\ulsbs@langsectfulltitle{\ulsbs@langcodewithtitle{\ulsbs@langsectlang}{\ulsbs@langsecttitle}}%
    \IfBooleanF{#1}{% if title is empty and there is no star, display no title
      \ifx\ulsbs@langsecttitle\empty%
        \def\ulsbs@langsectfulltitle{}%
      \fi%
    }%
    \IfStrEq{\ulsbs@langsectlang}{PT}{%
      \color{ulLanguagePortugueseColor}%
      \begin{portuguese}%
    }{}%
    \IfStrEq{\ulsbs@langsectlang}{EN}{%
      \color{ulLanguageEnglishColor}%
      \begin{english}%
    }{}%
    \IfStrEq{\ulsbs@langsectlang}{FI}{%
      \color{ulLanguageFinnishColor}%
      \begin{finnish}%
    }{}%
  }{% there was no lang set
    \def\ulsbs@langsectlang{}%
    \ifx\ulsbs@langsecttitle\empty%
      \def\ulsbs@langsectfulltitle{}%
    \else%
      \def\ulsbs@langsectfulltitle{\ulsbs@langsecttitle}%
    \fi%
  }%
  \ulsbs@langsectfont%
  % Restore font leading even when the font hook selects the size already active.
  \baselineskip=\normalbaselineskip%
  \ifhmode\par\fi%
  \ifx\ulsbs@langsectfulltitle\empty\else%
    \ifulsbs@upcaselangsecttitle%
      \ulsbs@langsectstart{\protect\MakeUppercase{\expandafter\ulsbs@langsectfulltitle}}% Create section
    \else%
      \ulsbs@langsectstart{\ulsbs@langsectfulltitle}% Create section
    \fi%
  \fi%
  \baselineskip=\ulsbs@langsectbaselinemultiplier\baselineskip plus 0pt minus 0pt%
  \IfBooleanTF{#4}{% 2nd star: obey lines
    \vspace{.25cm}%
    \obeylines%
  }{% no star: do not obey lines
    \disobeylines% use the renewed version above to set up parskip etc
  }%
  \advance\leftskip\ulsbs@langsectleftskipadd%
  \rightskip=\ulsbs@langsectrightskip%
  \parindent=0pt\relax%
  \noindent%
  \ulsbs@langsectbeginextra%
  \ignorespaces%
}{% post env
  \par%
  \IfStrEq{\ulsbs@langsectlang}{PT}{%
    \end{portuguese}%
  }{}%
  \IfStrEq{\ulsbs@langsectlang}{EN}{%
    \end{english}%
  }{}%
  \IfStrEq{\ulsbs@langsectlang}{FI}{%
    \end{finnish}%
  }{}%
  \ulsbs@langsectendextra%
  \egroup%
}

% === Passages ===

% Consumer defaults: define providecommand macros before loading or renew after
% loading; they are read when the passage starts. These lengths are allocated
% here, with 2em evaluated in the current post-typography font context. Change
% them after loading with \setlength, not \renewcommand; do not move their
% initialization into the pre-typography book-defaults phase.
\providecommand{\ulPassageFont}{\ulLyricFontStyle}%
% Style for the default text for alt. passage env (starred):
\providecommand{\ulAltPassageFont}{\ulLyricFontStyle\normalsize\sffamily\itshape}%
\providecommand{\ulPassageTitleFont}{\ulSongTitleFont}%
% Baseline multiplier, must be a number:
\providecommand{\ulPassageBaselineMultiplier}{1}%
% Amount of indentation for default text:
\newlength\ulPassageLeftSkipAdd\setlength\ulPassageLeftSkipAdd{2em}\relax%
\newlength\ulPassageRightSkip\setlength\ulPassageRightSkip{2em}\relax%
% The content in languages that have their code in in this comma separated
% list are not shown at all:
\providecommand{\ulPassageLanguagesToSkip}{}%
% The language *codes* held in this comma separated list are not displayed:
\providecommand{\ulPassageOmitLangCodes}{}%

% Starts a new subsubsection and prints it's title given as the only parameter.
% Used in the definition of passage environment.
\newcommand\ulsbs@passagestart{\@startsection{subsubsection}{3}{\ulContentIndentWidth}%\z@
                          {-3.25ex\@plus -1ex \@minus -.2ex}%
                          {1.3ex \@plus .2ex}%
                          {\ulPassageTitleFont}}

% Environment for a passage of text that can be put inside a song environment
% and elsewhere. Supports colouring and hyphenation of language texts.
% Uses ulsbs@langsect environment.
% Requires packages: xstring, xparse, xcolor
% Usage: \begin{passage}*[LANGCODE]{<title>}
%   - If star is present, another version of the environment is used, meant
%     for notes.
%   - If LANGCODE is present, it is displayed before the title in brackets.
%     If it is one of the supported ones, the passage text is coloured and
%     hyphenation for that language is used.
%   - If <title> is empty, no title is displayed. To display it in such a
%     case, make the title a whitespace: { }
%   - \obeylines and \disobeylines (the default) macros are supported within
% Policy and language hiding retain verbatim's raw line-oriented source skip;
% the terminator must be on its own source line, not inside a captured argument.
\NewDocumentEnvironment{passage}{s o m}{% pre env
  % Check if this environment needs hiding:
  \newif\ifulsbs@hide\ulsbs@hidefalse%
  \ifulShowPassage\else\ulsbs@hidetrue\fi%
  % Test for skipped languages:
  \IfValueT{#2}{%
    \IfSubStr{\ulPassageLanguagesToSkip}{#2}{%
      \ulsbs@hidetrue% environment-local decision, never a change to feature policy
    }{}%
  }%
  % If hiding is needed, hide the rest using verbatim's \comment:
  \ifulsbs@hide\expandafter\comment\fi%
  \bgroup%
  \renewcommand{\ulsbs@langsectstart}{\ulsbs@passagestart}%
  \renewcommand{\ulsbs@langsectfont}{\ulPassageFont}%
  \renewcommand{\ulsbs@langsecttitlefont}{\ulPassageTitleFont}%
  \renewcommand{\ulsbs@langsectbaselinemultiplier}{\ulPassageBaselineMultiplier}%
  \renewcommand{\ulsbs@langsectbeginextra}{}%
  \renewcommand{\ulsbs@langsectendextra}{\vfil}%
  \renewcommand{\ulsbs@langsectomitlangcodes}{\ulPassageOmitLangCodes}%
  \ulsbs@langsectleftskipadd=\ulPassageLeftSkipAdd%
  \ulsbs@langsectrightskip=\ulPassageRightSkip%
  \IfBooleanT{#1}{% starred version
    \renewcommand{\ulsbs@langsectfont}{\ulAltPassageFont}%
    \renewcommand{\ulsbs@langsectbeginextra}{\vspace{1cm}\vfil}%
  }
  \IfValueTF{#2}{%
    \begin{ulsbs@langsect}[#2]{#3}
  }{% no lang
    \begin{ulsbs@langsect}{#3}
  }%
}{% post env
  \ifulsbs@hide% End hidden section (\endcomment from verbatim package):
    \expandafter\endcomment%
  \else% Execute the post command(s):
    \end{ulsbs@langsect}
    \egroup%
  \fi%
}

% === In-song blocks, translations and explanations ===

\newcounter{ulsbs@aftersongblocks}\setcounter{ulsbs@aftersongblocks}{0}% reset to 0 in endsong

% Layout owns fill state and the in-song block-start implementation.

% Define \blockstart to be able to use in-song block feature manually also
% elsewhere than in the environments defined below.
\newcommand{\blockstart}{\ulsbs@blockstart}

% Displays the parameter text centered in a large font. Use within songs, but outside
% verses. Intended to be used for short mantras for recitation.
\newcommand{\showmantra}[1]{%
  \ulEndSongVfill{true}%
  % Do not call \ulsbs@blockstart here, as this is not really an after-song block...
  \vspace{1.618034em}%
  \begin{center}%
    {\Large #1}%
  \end{center}%
}

% Feeler prose, in songs or standalone intersong blocks. Capture balanced bodies
% without expansion; unlike environ's cleanup stream, lazy selection keeps hidden
% conditional tokens opaque. This is not the raw line-oriented prose skip.
\ExplSyntaxOn
\NewDocumentEnvironment {feeler} {+b}
  {\legacy_if:nT {ulShowFeeler}
    {
      \ulEndSongVfill{true}
      \ulsbs@blockstart
      \small
      \par\noindent\ignorespaces
      \begin{center}\begin{em}
        #1
      \end{em}\end{center}
    }}
  {\ignorespacesafterend}
\ExplSyntaxOff

% Translation font default: define before loading or renew afterward. It is
% applied inside each translation block, after the built-in \small selection.
\providecommand{\ulTranslationFont}{}
% Numeric multiplier applied to the leading selected by the translation font.
\providecommand{\ulTranslationBaselineMultiplier}{1}

\newif\ifulsbs@intranslation\ulsbs@intranslationfalse%

% Environment for song translation, to be used only inside a song.
% Keep verbatim's line-oriented skip: hidden source need not have balanced TeX
% braces/conditionals. The terminator must remain on its own source line.
\NewDocumentEnvironment{translation}{o}{% pre env
  % Check if this environment needs hiding
  \newif\ifulsbs@hide\ulsbs@hidefalse%
  \ifulShowTranslation\else\ulsbs@hidetrue\fi%
  % If hiding is needed, hide the rest using verbatim's \comment
  \ifulsbs@hide\expandafter\comment\fi%
  % The following, including the environment body, is only expanded and
  % processed, if no hiding is needed.
  \ulsbs@intranslationtrue%
  \ulEndSongVfill{true}%
  \newcommand{\nextverse}{\ulBreakPenalty=\ulTranslationVersePenalty\brk\vspace{.5em}}% some vertical space between verses
  \IfValueTF{#1}{% optional lang id (must be upcase) signifying translation
    \def\ulsbs@langsectlang{#1}%
    \IfStrEq{\ulsbs@langsectlang}{PT}{%
      \color{ulLanguagePortugueseColor}%
      \begin{portuguese}%
    }{}%
    \IfStrEq{\ulsbs@langsectlang}{EN}{%
      \color{ulLanguageEnglishColor}%
      \begin{english}%
    }{}%
    \IfStrEq{\ulsbs@langsectlang}{FI}{%
      \color{ulLanguageFinnishColor}%
      \begin{finnish}%
    }{}%
  }{% no lang id given
    \def\ulsbs@langsectlang{}%
  }%
  \beginverse*%
  \vspace{-\ulsbs@verse@separation}% remove the vertical space created by \beginverse
  \ulsbs@blockstart%
  \chordsoff%
  \small\ulTranslationFont%
  % Keep any elastic leading installed by an existing font hook at factor 1.
  \ifdim\ulTranslationBaselineMultiplier pt=1pt\relax\else
    \baselineskip=\ulTranslationBaselineMultiplier\baselineskip plus 0pt minus 0pt\relax
  \fi%
}{% post env
  \ifulsbs@hide% End hidden section (\endcomment from verbatim package):
    \expandafter\endcomment%
  \else% Execute the post command(s):
    \endverse%
    \IfStrEq{\ulsbs@langsectlang}{PT}{%
      \end{portuguese}%
    }{}%
    \IfStrEq{\ulsbs@langsectlang}{EN}{%
      \end{english}%
    }{}%
    \IfStrEq{\ulsbs@langsectlang}{FI}{%
      \end{finnish}%
    }{}%
  \fi%
}

% Explanation defaults: define before loading or renew afterward. Font hooks,
% numeric baseline multiplier and language lists are read at environment entry.
\providecommand{\ulExplanationFont}{\rmfamily\footnotesize}%
\providecommand{\ulExplanationTitleFont}{\bfseries}%
% Baseline multiplier, must be a number:
\providecommand{\ulExplanationBaselineMultiplier}{1}%
% The content in languages that have their code in in this comma separated
% list are not shown at all:
\providecommand{\ulExplanationLanguagesToSkip}{}%
% The language *codes* held in this comma separated list are not displayed:
\providecommand{\ulExplanationOmitLangCodes}{}%

\newcounter{ulsbs@explcounter}\setcounter{ulsbs@explcounter}{0} % reset to 0 in endsong

% Environment for 'song explanation'
% Uses ulsbs@langsect environment
% Requires packages: xstring, xparse, xcolor
% Usage: \begin{explanation}[LANGCODE]
%   - If LANGCODE is present, it is displayed before the title in brackets.
%     If it is one of the supported ones, the passage text is coloured and
%     hyphenation for that language is used.
% Policy and language hiding retain verbatim's raw line-oriented source skip;
% the terminator must be on its own source line, not inside a captured argument.
\NewDocumentEnvironment{explanation}{o}{% pre env
  % Check if this environment needs hiding:
  \newif\ifulsbs@hide\ulsbs@hidefalse%
  \ifulShowExplanation\else\ulsbs@hidetrue\fi%
  % Test for skipped languages:
  \IfValueT{#1}{%
    \IfSubStr{\ulExplanationLanguagesToSkip}{#1}{%
      \ulsbs@hidetrue% environment-local decision, never a change to feature policy
    }{}%
  }%
  % If hiding is needed, hide the rest using verbatim's \comment:
  \ifulsbs@hide\expandafter\comment\fi%
  \bgroup%
  \ifnum\value{ulsbs@explcounter}<1%
    \ulsbs@blockstart%
  \else%
    \vspace{1ex}%
  \fi%
  \addtocounter{ulsbs@explcounter}{1}%
  \ulEndSongVfill{false}% so far explanation is last block, no vfill needed at song end
  \setlist[description]{% modify description list (enumitem package)
    %topsep=30pt,%      % space before start / after end of list
    %itemsep=5pt,%      % space between items
    leftmargin=3em%     % left margin
    %font={\rmfamily}%   % set the label font
  }%
  \ulsbs@langsectleftskipadd=0pt\relax%
  \ulsbs@langsectrightskip=0pt\relax%
  \renewcommand{\ulsbs@langsectbaselinemultiplier}{\ulExplanationBaselineMultiplier}%
  \renewcommand{\ulsbs@langsectfont}{\ulExplanationFont}%
  \renewcommand{\ulsbs@langsecttitlefont}{\ulExplanationTitleFont}%
  \renewcommand{\ulsbs@langsectendextra}{}%
  \renewcommand{\ulsbs@langsectomitlangcodes}{\ulExplanationOmitLangCodes}%
  \IfValueTF{#1}{% language code is given
    \renewcommand\ulsbs@langsectstart[1]{\noindent{\ulsbs@langcodeformatted{#1}}}%
    \begin{ulsbs@langsect}*[#1]{}%
  }{% no language code
    \begin{ulsbs@langsect}{}%
  }%
}{% post env
  \ifulsbs@hide% End hidden section (\endcomment from verbatim package):
    \expandafter\endcomment%
  \else% Execute the post command(s):
    \end{ulsbs@langsect}%
    \egroup%
  \fi%
}
